\subsection{The completion of a field with a valuation}

PROPOSITION 5. Let K be a not necessarily commutative field, v a valuation on K and G the group \(v(\mathbf{K}^{*})\) with the discrete topology.

\begin{enumerate}
\def\labelenumi{(\alph{enumi})}
\item
  The completion ring \(\hat{\mathbf{K}}\) of \(\mathbf{K}\) (with \(\mathcal{T}_v\) ) is a topological field.
\item
  The mapping \(v: \mathbf{K}^* \to \mathbf{G}\) can be extended uniquely to a continuous mapping \(\hat{v}: \hat{\mathbf{K}}^* \to \mathbf{G}\) . The mapping \(\hat{v}\) (extended by \(\hat{v}(0) = +\infty\) ) is a valuation on \(\hat{\mathbf{K}}\) and \(\hat{v}(\hat{\mathbf{K}}^*) = v(\mathbf{K}^*)\) .
\item
  The topology on \(\hat{\mathbf{K}}\) is the topology defined by the valuation \(\hat{v}\) .
\item
  For all \(\alpha \in \mathbf{G}\) let \(\mathrm{V}_{\alpha}, \mathrm{V}_{\alpha}'\) be the subgroups of \(\mathbf{K}\) defined by the conditions \(v(x) > \alpha\) , \(v(x) \geqslant \alpha\) . Then the closures \(\overline{\mathrm{V}}_{\alpha}, \overline{\mathrm{V}}_{\alpha}'\) of \(\mathrm{V}_{\alpha}, \mathrm{V}_{\alpha}'\) in \(\hat{\mathbf{K}}\) are defined by the conditions \(\hat{v}(x) > \alpha\) , \(\hat{v}(x) \geqslant \alpha\) respectively.
\item
  The ring of \(\hat{v}\) is the completion \(\hat{\mathbf{A}}\) of the ring \(\mathbf{A}\) of \(v\) ; the ideal of \(\hat{v}\) is the completion \(\hat{\mathfrak{m}}\) of the ideal \(\mathfrak{m}\) of \(v\) .
\item
  \(\hat{\mathbf{A}} = \mathbf{A} + \hat{\mathfrak{m}}\) ; the residue field of \(\hat{v}\) is canonically identified with that of \(v\) .
\end{enumerate}

To prove (a) it suffices (General Topology, Chapter III, § 6, no. 8, Proposition 7) to show the following: let \(\mathfrak{F}\) be a Cauchy filter (with respect to the additive uniform structure) on \(\mathbf{K}^*\) for which 0 is not a cluster point; then the image of \(\mathfrak{F}\) under the bijection \(x \mapsto x^{-1}\) is a Cauchy filter (with respect to the additive uniform structure). For since 0 is not a cluster point of \(\mathfrak{F}\) , there exists \(\mathbf{M} \in \mathfrak{F}\) and \(\beta \in \mathbf{G}\) such that \(\beta\) is an upper bound of \(v(\mathbf{M})\) . Let \(\alpha \in \mathbf{G}\) . If \(\mathbf{M}'\) is an element of \(\mathfrak{F}\) contained in \(\mathbf{M}\) and such that \(v(x - y) > \sup (\alpha + 2\beta, \beta)\) for \(x \in \mathbf{M}'\) and \(y \in \mathbf{M}'\) , then \(v(x^{-1} - y^{-1}) > \alpha\) for \(x \in \mathbf{M}'\) and \(y \in \mathbf{M}'\) (no. 1, Lemma 1). Whence (a).

By Proposition 1 of no. 1, \(v|K^*\) is a continuous homomorphism from \(K^*\) to \(G\) and hence can be extended uniquely to a continuous homomorphism \(\hat{v}\) from \(K^*\) to \(G\) . the Relation

\[
\hat {v} (x + y) \geqslant \inf (\hat {v} (x), \hat {v} (y))
\]

holds in \(\mathbf{K}^*\) and hence also holds in \(\hat{\mathbf{K}}^*\) by continuity. Thus \(\vartheta\) (extended by \(\hat{\upsilon}(0) = +\infty\) ) is a valuation on \(\hat{\mathbf{K}}\) and (b) is proved.

We now show (d). Let \(\alpha \in G\) and \(x \in \overline{V}_{\alpha} - \{0\}\) . For \(y\) in \(V_{\alpha}\) sufficiently close to \(x\) , \(\hat{v}(x) = \hat{v}(y) = v(y)\) and hence \(\hat{v}(x) > \alpha\) . Conversely, let \(x \in \widehat{K}^*\) be such that \(\hat{v}(x) > \alpha\) ; for \(y\) in \(K^*\) sufficiently close to \(x\) , \(v(y) = \hat{v}(y) = \hat{v}(x)\) and therefore \(y \in V_{\alpha}\) , whence \(x \in \overline{V}_{\alpha}\) . Thus \(\overline{V}_{\alpha}\) is the set of \(x \in \widehat{K}\) such that \(\hat{v}(x) > \alpha\) . The argument is analogous for \(V_{\alpha}'\) . This proves (d).

Taking account of Proposition 7 of General Topology, Chapter III, § 3, no. 4, assertion (c) is a consequence of (d). Assertion (e) is a special case of (d). Finally let \(x \in \hat{\mathbf{A}}\) ; there exists \(y \in \mathbf{A}\) such that \(\hat{v}(x - y) > 0\) ; then \(z = x - y \in \hat{\mathfrak{m}}\) and hence \(x = y + z \in \mathbf{A} + \hat{\mathfrak{m}}\) ; thus \(\hat{\mathbf{A}} = \mathbf{A} + \hat{\mathfrak{m}}\) , which shows (f).

Remark. For all \(x \in \hat{\mathbf{K}}\) not belonging to \(\hat{\mathbf{A}}\) , there exists \(x_0 \in \mathbf{K}\) such that \(\hat{v}(x - x_0) > 0\) , \(\hat{v}(x) = \hat{v}(x_0) = v(x_0) < 0\) ; then \(x_0^{-1}x \in \hat{\mathbf{A}}\) and, as \(x_0^{-1} \in \mathbf{A}\) , it is seen that, if we set \(S = A - \{0\}\) , it is possible to write \(\hat{\mathbf{K}} = S^{-1}\hat{\mathbf{A}}\) .

